% Paper 3 — Introduction (borrador 2026-10-17). Números: docs/paper3-outline-v3-2026-10-17.md
% (bloque "Números congelados") y los JSON docs/paper3-*.json generados sobre
% docs/paper3-archive-manifest.txt.

\section{Introduction}\label{sec:intro}

The harness around a language model (its prompts, tools, context policy and
control loop) has become an object of optimization in its own right. Since early
2026, systems such as Meta-Harness, AutoDesign and Harness Continual Learning
let a proposer agent rewrite the harness while the model stays fixed, and keep a
rewrite when it scores better than the incumbent
\citep{lee2026metaharness,luo2026autodesign,kang2026hcl}. Every such loop ends
in the same decision: commit the candidate or discard it. Because agent
evaluations are stochastic, that decision is a statistical test whether or not
it is written as one, and its error rate determines whether the loop climbs or
drifts.

Two recent observations make the question pressing. First, the gains of harness evolution are fragile: they are unstable and transfer only partially to held-out tasks \citep{wu2026harnessdev}, and on Terminal-Bench 2.1, under matched budgets, evolution does not beat simple test-time scaling, with long-horizon games as the exception \citep{wang2026rethinking}. Second, the noise is large where these loops are
cheapest to run. On our own archive of agent sweeps with small local models,
the median configuration flips 5.7\% of its tasks between two runs of the
\emph{same} setup, and one configuration in ten flips more than 17\%. A gate
that commits on a point estimate is answering a coin toss on a sizeable fraction
of its decisions.

The field has started to respond. Between June and September 2026, four gates
replaced the point estimate with a statistical test: an anytime-valid
e-process \citep{zayxshawn2026pace}, a paired $z$-test
\citep{luo2026harnessbank}, a one-sided paired test on discordant items
\citep{wang2026sage}, and a noise tolerance set from repeated evaluations of the
base harness \citep{xia2026rrsi}. The last one is the natural reference point
for this paper, because it calibrates the threshold against measured noise.
It measures that noise once, on the base harness, and keeps the tolerance fixed
for the rest of the evolution.

This paper asks what that single calibration costs. The answer depends on a fact
that, to our knowledge, has not been measured: whether the run-to-run noise of
an agent configuration is a property of the benchmark or of the configuration.
If it is a property of the benchmark, one measurement on the base suffices. If
it changes with the configuration, a tolerance measured on the base is a bet on
the noise of incumbents that have not been measured yet.

We answer it with an archive of 140 sweeps that a single project produced over
four months (32 backend specifications, mostly small models served locally,
agentic coding suites graded by compilation and assertions, and 14{,}790 runs in
groups that repeat an identical configuration). Because every pair of runs of an identical configuration is a
null comparison, any change a gate accepts on such a pair is a false acceptance,
and we can count those acceptances without simulating effects. We make four
contributions:

\begin{enumerate}
  \item \textbf{Noise is a property of the configuration.} Within a single sweep
  on a single suite, the flip rate of identical runs differs across
  configurations by a median of 11.1 percentage points (up to 26.3), and 29 of
  90 incumbents are noisier than the base of their sweep
  (Section~\ref{sec:res-heterogeneity}).
  \item \textbf{A threshold calibrated once does not control false acceptances.}
  With the tolerance set to the spread of the base's repeated runs, as in
  \citet{xia2026rrsi}, the gate accepts 6.9\% of null changes across 1{,}814
  null pairs; re-deriving the tolerance from each configuration's own repeated
  runs lowers this to 0.8\%. The single calibration fails worst when the base
  happens to be quiet (17.8\%) and is safe only when the base happens to be the
  noisiest configuration in the sweep (Section~\ref{sec:res-fixed}).
  \item \textbf{Re-deriving the threshold is not a trade-off.} Injecting real
  improvements of known size, the per-configuration rule matches or exceeds the
  power of the fixed rule at every effect size, because a fixed tolerance is
  too narrow for noisier incumbents and too wide for quieter ones. Its cost is
  runs, not power: each incumbent must be repeated (Section~\ref{sec:res-power}).
  \item \textbf{The threshold curve.} Across 1{,}614 null comparisons, the
  decision rules of published optimizers accept between 1.1\% and 26.3\% of null
  changes, and the same rule moves by more than an order of magnitude when only
  its margin changes. No gate we are aware of reports this curve
  (Section~\ref{sec:res-curve}).
\end{enumerate}

The scope of these claims is deliberate. Our regime is mostly small local models with
large capability gaps, where noise is high and heterogeneous; the frontier models
that most self-evolving systems optimize may show less heterogeneity, and we do
not claim that their published improvements are false. What we claim is
narrower: wherever the noise of configurations differs, a threshold calibrated
once is miscalibrated for most of the decisions it makes, and the error is
measurable from data that any repeated-run protocol already produces.
